A bank that sets the PD of a grade from the published statistics of rating agencies inherits two choices it seldom shows: which events count as a default, and what becomes of a rating withdrawn during the year. Estimators take the published tables as given, and the European repository that publishes them sets no common definition of a default event. We propose a measurement protocol that treats the definition and the treatment of withdrawals as parameters of the PD: the PD of every grade under each page of the repository with its capital, the bias of each treatment of withdrawn ratings on rating paths of known truth, matrix-test rejection rates, and the exact coverage of interval estimators. On S\&P's corporate cohorts of 2000 to 2025 the long-run PD of CCC to C is $\spLraDtwoCcc\%$ on the default page and $\spLraDfourCcc\%$ on the transition page, and the IRB risk weight of the cohort mix falls from $\spRwDtwo\%$ to $\spRwDfour\%$. With withdrawal ten times as likely in the year before default, removing the withdrawn ratings understates the PD of grade B by $\removedBkNine\pm\removedBkNinese$ percent and following them to the year end by $\followedBkNine\pm\followedBkNinese$ percent, while two matrix tests reject in $\testthkNinefive\%$ and $\testmrkNinefive\%$ of repetitions at the 5\% level. Closed forms for the bias of each treatment agree with the simulation within $\thMaxZ$ standard errors. With 50 obligors per grade the Wald interval contains the PD of AAA in $\waldAaaFifty\%$ of samples. The protocol measures these biases and repairs none.
